\part{Álgebra}
\chapter{Básicos}
Aquí exploraremos los conceptos básicos
relacionados al álgebra. El álgebra
aparece en todas partes de la matemática,
y es importante conocer lo básico
ya que en particular se puede aplicar a
infinidad de problemas.


\section{Inducción}
Esto es un método fundamental en todas
las áreas de la matemática.

Para ilustrar qué dice lo podemos imaginar como
una serie de dominós en fila, tal que
si cae uno de ellos entonces cae el que
está inmediatamente a la derecha.
Es claro que si cae el primer dominó, entonces
caen todos los demás ya que si cae el primero,
entonces cae el segundo; si cae el segundo
entonces cae el tercero, y así sucesivamente.

% TODO: Añadir figura de dominós cayendo

Esta misma idea se puede formalizar, pero
en lugar de dominós, con proposiciones.
Esta es una definición informal:

\begin{theorem}[Inducción]
	\label{theorem:induccion}
	Si $P(n)$ es una proposición cuya veracidad
	depende del valor de $n$, se sabe que si
	$P(n)$ es cierto entonces $P(n + 1)$ es
	cierto, y $P(1)$ es cierto, entonces $P(n)$
	es cierto para todo $n$.
\end{theorem}

Para digerir el método lo mejor es usar
ejemplos.

\begin{example}[Suma Gaussiana]
	Demuestre que la suma de los primeros
	$n$ números naturales es igual a
	\[\frac{n(n + 1)}2.\]
\end{example}

Este resultado se le atribuye
al matemático Carl Gauss, quien, se cuenta,
lo dedujo cuando era un niño.

La manera de imaginar por qué
este resultado es cierto es sumar
los números de extremo a extremo como se
ve en la figura:

\begin{figure}[H]
	\begin{center}
		\includegraphics{figures/gaussian_sum.pdf}
	\end{center}
	\caption{Suma Gaussiana.}
\end{figure}

Por el momento simplificamos el problema
y consideramos que $n$ es par. Claramente
la suma del primer con el último
número es $n + 1$,
del segundo con el penúltimo $2 + n - 1 = n + 1$,
y así sucesivamente todos los pares tienen la
misma suma. Pero nos damos cuenta de que
hay exactamente $\frac n2$ pares, de ahí la suma
necesariamente debe ser igual a
\[\frac n2 \cdot (n + 1).\]

Para demostrar con inducción que la fórmula
es cierta para todo $n$
podemos verificar que la misma funciona
para $n = 1$ ya que si sustituimos obtenemos
\[\frac{1 \cdot 2}{2} = 1\]
que es la suma hasta el primer número natural.

Ahora notemos que si conseguimos probar que
si la fórmula funciona para $n$, entonces
funciona para $n + 1$, entonces funcionaría
para todo $n$ ya que la misma funciona para
$n = 1$\footnote{Ten en cuenta el \autoref{theorem:induccion}}

Para ello si ocurre que
\[1 + 2 + \cdots + n = \frac{n(n + 1)}2,\]
debemos probar que
\[1 + 2 + \cdots + n + (n + 1) = \frac{(n + 1)(n + 2)}2,\]
pero sustituyendo
\[1 + 2 + \cdots + n = \frac{n(n + 1)}2\]
notamos que
\[\frac{n(n + 1)}2 + (n + 1) = (n + 1)\left(\frac n2 + 1\right) = \frac{(n + 1)(n + 2)}2,\]
que es lo que queríamos probar.

\subsection*{Problemas para esta sección}
\begin{problem}
Muestre que
\[1^2 + 3^2 + \cdots + n^2 = \frac{n(n + 1)(2n + 1)}6.\]
\end{problem}

\begin{problem}
Pruebe que la suma de los primeros
$n$ números impares es igual a $n^2$.
\end{problem}

\section{Identidades}
Manejar expresiones algebraicas puede ser
bastante complicado si no se es capaz
de simplificar la misma.

E aquí algunos métodos y sustituciones
útiles.

\subsection{Trinomio cuadrado perfecto}
El trinomio cuadrado perfecto es quizás
el caso de factoreo más común de todos,
y tan solo consiste en saber que
\[(a + b)^2 = a^2 + 2ab + b^2.\]

En los problemas, a veces es conveniente
expandir la expresión $(a + b)^2$
y en otros casos resulta útil simplificar
$a^2 + 2ab + b^2$ como $(a + b)^2$.

\subsection{Super identidad}
Un amigo me contó que esta simple
transformación aparece en muchos lugares
inesperados y especialmente en teoría de números:
\[ab + cd = a(b - c) + c(d + a)\]

\subsection{Identidad de Sophie Germain}
La identidad de Sophie Germain establece lo siguiente:
\[a^4 + 4b^4 = (a^2 + 2b^2 + 2ab)(a^2 + 2b^2 - 2ab)\]

\subsection{Racionalización}
Normalmente es difícil tratar con raíces en
los denominadores de fracciones,
y el simple hecho de que
\[(a + b)(a - b) = a^2 - b^2\]
nos ayuda bastante ya que si
$a$ o $b$ son raíces cuadradas, es posible
multiplicar la expresión por $a - b$
para deshacerse de las raíces,
así, un uso común es cuando tenemos algo del
siguiente estilo:

\[\frac1{\sqrt a + \sqrt b}.\]

Podemos multiplicar la fracción
por $1$ y eso no cambia su valor,
pero
\[1 = \frac{\sqrt a - \sqrt b}{\sqrt a - \sqrt b}.\]

y si lo aplicamos a la expresión anterior vemos que
\[\frac1{\sqrt a + \sqrt b}\cdot\frac{\sqrt a - \sqrt b}{\sqrt a - \sqrt b} = \frac{\sqrt a - \sqrt b}{\sqrt a ^2 - \sqrt b ^2} = \frac{\sqrt a - \sqrt b}{a - b}.\]

Y ahora ya no tenemos raíces en el denominador.
Esta transformación es especialmente
últil si $a$ y $b$ son números
enteros o racionales, o
si $a - b$ lo es.

\chapter{Funciones}
Si ya sabes qué es una función,
entonces puedes saltar este capítulo.

\section{Fundamentos}

Una función puede ser pensada como
una máquina que toma un valor o una serie
de valores y devuelve un resultado.
Por ejemplo, si recibe como entrada
$x$ y produce una salida $y$, eso
se puede pensar así:

\begin{figure}[H]
	\centering
	\includegraphics{figures/maquinafuncion.pdf}
	\caption{Imagen mental de una función.}
\end{figure}

Por ejemplo, si esta máquina se llama $f$
y si como entrada le asignamos el número
$1$ y $f$ devuelve el número $2$, entonces
eso se simboliza así:
\[f(1) = 2.\]
En general, el valor $f(x)$ simboliza la
salida que devuelve $x$ si le asignamos $x$
como la entrada.

Como se puede ver, este concepto es
muy general y puede ser aplicado a
casi cualquier cosa.

Hablando de una manera más rigurosa,
debemos definir qué es exactamente una función,
pero antes necesitamos otro concepto.

\begin{definition}[Producto cartesiano]
	Sea $A$ y $B$ dos conjuntos,
	el \textit{Producto Cartesiano} de estos
	dos conjuntos es definido por
	\[A \times B \coloneqq \{(a, b) : a \in A, b \in B\}.\]
\end{definition}

La interpretación de esto es que $A \times B$
es un conjunto de pares ordenados donde
el primer elemento del par ordenado
es un elemento de $A$ y el segundo un elemento
de $B$. Por ejemplo, si $A = \{1, 2, 3\}$
y $B = \{3, 4\}$, entonces los elementos
de $A \times B$ son
\[
	\begin{array}{cccccc}
		(1, 3) & (1, 4) & (2, 3) & (2, 4) & (3, 3) & (3, 4)
	\end{array}
\]

También podemos definir el producto cartesiano
entre otros conjuntos, como por ejemplo
$[1, 2] \times [2, 3]$, y se puede visualizar de la
siguiente manera:
\begin{figure}[H]
	\centering
	\includegraphics{figures/cartesianproduct.pdf}
	\caption{El producto cartesiano.}
\end{figure}

\begin{definition}[Función]
	Una función $f$ que `va' de un conjunto
	$A$ a un conjunto $B$ es
	un subconjunto de $A \times B$ tal que
	para todo $a \in A$, existe un único
	elemento $b \in B$ con $(a, b) \in f$.
	Esto se simboliza con
	\[f \colon A \to B.\]
\end{definition}

Al primer conjunto $A$ se le llama
\textbf{dominio} de la función $f$
y se denota por $D(f)$,
y al conjunto $B$ el \textbf{codominio}.
El conjunto de todos los valores
de salida de la función se llama
el \textbf{rango} y se denota
por $R(f)$ (notemos que $R(f) \subseteq B$),
así
\[R(f) = \{b \in B \colon \exists a \in A : f(a) = b\}.\]

Otra manera de interpretar que un subconjunto
de $A \times B$ es una función es imaginar
ese conjunto en el plano cartesiano
y verificar que ninguna recta \textbf{vertical} corta
al conjunto en más de un punto.

\begin{figure}[H]
	\centering
	\includegraphics{figures/funciones.pdf}
\end{figure}

O con conjuntos:

\begin{center}
	\includegraphics{figures/conjuntosfunciones.pdf}
\end{center}

\begin{definition}[Inyectividad]
	Una función es \textit{inyectiva} si
	\[f(a) = f(b) \iff a = b.\]
\end{definition}

Esto se interpreta conque la función $f$
no asigna dos valores del dominio
al mismo valor del rango, o sea
que no puede darse $f(a) = f(b)$
con $a \neq b$.

También puede verse como que ninguna
recta \textbf{horizontal} del plano corta a
a la función en más de un punto.

\begin{figure}[H]
	\centering
	\includegraphics{figures/inyectividad.pdf}
\end{figure}

\begin{definition}[Sobreyectividad]
	Una función $f\colon A \to B$ es sobreyectiva
	si \textbf{para todo} $b \in B$, existe
	por lo menos
	un $a \in A$ tal que $f(a) = b$.
\end{definition}

Esto quiere decir que la función adquiere
todos los valores de su codominio.

\begin{figure}[H]
	\centering
	\includegraphics{figures/funcionsobreyectiva.pdf}
	\caption{Función sobreyectiva.}
\end{figure}

\section{Definiendo funciones}
Para definir una función, solo necesitamos
dos conjuntos no vacíos, y por simpleza consideramos
los conjuntos $A = B = \{1, 2, 3\}$.
Claramente podemos definir una función $f \colon A \to B$
de esta manera:
\[f = \{(1, 2), (2, 2), (3, 1)\}.\]
Así, tendremos
\[
	\begin{array}{ccc}
		f(1) = 2 & f(2) = 2 & f(3) = 1.
	\end{array}
\]

Notemos que el valor $3$ del codominio
no está en el rango de $f$, al contrario de $1$ y $2$.

Si definimos el conjunto
\[g = \{(1, 2), (1, 3), (2, 3), (3, 1)\},\]
entonces $g$ no es una función,
porque los pares ordenados
$(1, 2)$ y $(1, 3)$ asignan el valor $1$
del dominio a más de un valor del codominio.

De manera similar, si
\[h = \{(1, 2), (3, 3)\},\]
entonces $h$ no sería una función de
$A$ a $B$ ya que el valor $2$ de
$A$ no es asignado a ningún elemento
de $B$.
\begin{remark}
	Nótese que, en cambio, $h$ sí es una función
	de $A\backslash\{2\}$ a $B$ ya que en este caso
	todos los valores del dominio
	son asignados a un único valor del codominio.
\end{remark}

Podemos definir funciones del conjunto de
los números reales al mismo conjunto,
digamos $f \colon \RR \to \RR$, por una
fórmula que devuelve un resultado, un ejemplo sería
el siguiente

\begin{example}
	Sea $f \colon \RR \to \RR$ una función
	definida por
	\[f(x) = x^2.\]
\end{example}

Básicamente esta función asigna cada número real a
su cuadrado.

Si analizamos la función vemos que $f$ no es inyectiva,
porque $f(1) = f(-1) = 1$, o sea que hay dos valores del
dominio que son asignados al mismo valor del codominio
(en general, si $x = -x$, entonces $f(x) = f(-x)$).
Vemos que $f$ tampoco es sobreyectiva,
porque nunca adquiere valores negativos.

Si vemos atentamente, nos damos cuenta que
en realidad decir que $f \colon \RR \to \RR$
está ``definida por'' $f(x) = x^2$ es una abreviación para decir
\[f = \{(a, b^2) : a, b \in \RR\}.\]

\begin{example}[Función que no hace nada]
	Sea $f \colon \RR \to \RR$ definida por
	\[f(x) = x.\]
\end{example}

Esta función sí es inyectiva, y también sobreyectiva.
\begin{question*}
	¿Por qué?
\end{question*}

Pero notamos que, en sí, esta función asigna
cada elemento del codominio a un único
elemento del dominio \textbf{y viceversa},
a esta clase de funciones se las denomina
funciones \textbf{biyectivas}.

\begin{definition}[Biyectividad]
	Una función es \textit{biyectiva}
	si asigna cada elemento de su dominio
	un único elemento de su codominio
	y viceversa. Más específicamente,
	es biyectiva si es inyectiva y sobreyectiva.
\end{definition}

Para ilustrar por qué si se cumplen estas
dos condiciones entonces la función es biyectiva podemos
hacer el siguiente análisis.
Vemos que si $f$ es inyectiva, entonces
casi es biyectiva, porque no hay dos
elementos del dominio que van
al mismo elemento del codominio,
o sea que si recortamos los elementos del
codominio que no usamos, nos quedaría una asignación
uno-a-uno de los elementos del dominio y codominio.
En cambio, si la función es a su vez sobreyectiva, entonces
no es necesario hacer este recorte y ya tendríamos
una asignación uno-a-uno de los elementos del
dominio con los elementos del codominio.

\begin{figure}[H]
	\centering
	\includegraphics{figures/funcionbiyectiva.pdf}
	\caption{Función biyectiva}
\end{figure}

\begin{lemma}[Muy útil]
	Sea $f \colon A \to A$ una
	función. Demuestre que si
	\[f(f(x)) = x\]
	para todo $x \in A$, entonces $f$
	es biyectiva.
\end{lemma}

\begin{proof}
	Vemos que para todo $x \in A$ (codominio de $f$),
	entonces $a = f(x) \in A$ satisface $f(a) = x$,
	por lo que es obvio que $f$ es sobreyectiva.

	Por el contrario, si
	\begin{align*}
		f(a)                 & = f(b)         \\
		\implies a = f(f(a)) & = f(f(b))  = b \\
		\implies      a      & = b,
	\end{align*}
	por lo que $f$ es inyectiva, completando la prueba.
\end{proof}

\begin{example}
	Sea $f \colon \ZZ \to \ZZ$ definida por
	\begin{equation*}
		f(x) =
		\begin{cases}
			n + 1 \hspace{1cm} \text{si $n$ es par} \\
			n - 1 \hspace{1cm} \text{si $n$ es impar.}
		\end{cases}
	\end{equation*}
	Demuestre que $f$ es biyectiva.
\end{example}
Este problema presenta una función no
muy complicada, pero cuyo análisis
se trivializa con el lema.

\paragraph{Caso 1. $n$ es par. } Entonces
\[f(f(n)) = f(n + 1) = (n + 1) - 1 = n.\]

\paragraph{Caso 2. $n$ es impar. } Análogamente
\[f(f(n)) = f(n - 1) = (n - 1) + 1 = n.\]

Por el lema que vimos recién, entonces $f$
es biyectiva porque $f(f(n)) = n$ para todo
$n$, sin importar la paridad de $n$.

\subsection*{Problemas para esta sección}
\begin{problem}
Sean $A = \{1, 2, 3, 4\}$ y $B = \{5, 6, 7\}$,
determine si los siguientes conjuntos
son funciones de $A$ a $B$:
\begin{itemize}
	\ii $\{(1, 5), (3, 6), (2, 6)\}$.
	\ii $\{(1, 4), (2, 5), (3, 6), (4, 7)\}$.
	\ii $\{(2, 5), (1, 5), (3, 5), (4, 5)\}$.
	\ii $\{(2, 5), (3, 6), (2, 7), (1, 5)\}$.
	\ii $\{(1, 5), (3, 7), (2, 7), (4, 6)\}$.
\end{itemize}
\end{problem}

\begin{problem}
Muestre que si $A$ y $B$ son
conjuntos finitos no vacíos
\textbf{del mismo tamaño},
entonces cualquier función
inyectiva de $A$ a $B$ es una
biyección de $A$ a $B$.
\end{problem}

\begin{problem}
Sea $f \colon \RR\backslash\{0\} \to \RR$ definida por
\[f(x) = \frac1x.\]
Demuestre que $f$ es biyectiva.
\begin{hints}
	\begin{hint}
		Una manera elegante es usar el problema anterior.
	\end{hint}
\end{hints}
\end{problem}

\section{Desigualdades}

\section{Ecuaciones}

\section{Problemas de álgebra}

\begin{problem}
Resuelva el sistema de ecuaciones:
\begin{equation*}
	\begin{cases}
		a^2 + b^2 + ab = 7 \\
		a - b          = 2.
	\end{cases}
\end{equation*}
\end{problem}

\begin{problem}
Si $a$ es un número real distinto de $0$, encuentre el valor de la expresión:
\[\frac{a + a^2 + a^3 + \cdots + a^n}{\frac{1}{a} + \frac{1}{a^2} + \frac{1}{a^3} + \cdots + \frac{1}{a^n}}\]
\end{problem}

\begin{problem}[Omapa 2023 - Ronda Departamental]
Encuentre el valor de la expresión
\[\frac2{\sqrt1 + \sqrt2} + \frac1{\sqrt2 + \sqrt3} + \cdots + \frac1{\sqrt{2022} + \sqrt{2023}}.\]
\, %%%% <----- WARNING: Si quitás el espacio la compilación crashea
\begin{hints}

	\begin{hint}
		Racionalización.
	\end{hint}
\end{hints}
\end{problem}

\begin{problem}[Libro ruso]
Demuestre que si los números positivos $a$, $b$ y $c$
forman una progresión aritmética,
entonces los números
\[\frac1{\sqrt b + \sqrt c}, \frac1{\sqrt c + \sqrt a}, \frac1{\sqrt a + \sqrt b}\]
también forma una progresión aritmética.
\end{problem}

\begin{problem}
Encuentre el valor de la expresión:
\[\frac1{1\cdot2} + \frac1{2\cdot3} + \frac1{3\cdot4} + \cdots + \frac1{n(n + 1)}.\]

% \begin{hints}
% 	\begin{hint}
% 		$\frac{1}{n} - \frac{1}{n + 1} = \frac1{n(n + 1 )}$.
% 	\end{hint}
% \end{hints}
\end{problem}

\begin{problem}
Se eligen $n$ números de la siguiente
figura, uno por cada fila y columna:

\begin{figure}[H]
	\begin{center}
		\includegraphics{figures/gridconstant.pdf}
	\end{center}
\end{figure}

Encuentre todos los posibles valores
para la suma de todos los números seleccionados.
\end{problem}
